\section{Introduction}
\label{sec:intro}

A standard claim in the post-quantum-blockchain literature is that
GPU acceleration delivers an order-of-magnitude speedup for lattice
operations. We measured this on shipping hardware and find the claim
is half-true: it depends entirely on \emph{batch size}.

Lux runs a post-quantum threshold signature (Pulsar / LP-073) whose
hot path is the canonical Number-Theoretic Transform over the ring
$R_q = \mathbb{Z}_q[X]/(X^{256}+1)$ with $q = 0x1000000004A01$
(49 bits, NTT-friendly: $q - 1 = 2^{49} - \cdot$ divisible by $2N$
for $N = 256$). Every signature, every key share rotation, every
threshold ceremony reduces to thousands of these NTTs. The CPU
implementation is the canonical Lattigo-derived Montgomery body
(luxfi/math/ntt at LP-107 v1.4.0); the same kernel is realized as
a Metal shader and a WGSL shader, byte-equal to CPU by KAT
construction.

This paper reports the empirical per-polynomial throughput on
Apple M1 Max across the three backends as a function of batch size,
identifies the crossover where each GPU backend first beats CPU,
and derives three scaling laws that govern when GPU acceleration
actually pays off for blockchain workloads.

\paragraph{Summary of findings.}
On Apple M1 Max:
\begin{itemize}
\item CPU per-NTT cost is essentially flat at $\sim 1.4$ \textmu s
across batch sizes 1 to 65{,}536, indicating near-perfect SIMD
saturation of the underlying loop-unrolled-by-16 path.
\item Metal pays a fixed $\sim 1.5$ ms per command-buffer dispatch.
At $n_\text{polys} = 1$ this is $\sim 1{,}100\times$ slower than CPU;
at $n_\text{polys} = 16{,}384$ Metal first ties CPU
($1.19$ \textmu s vs $1.49$ \textmu s); at $n_\text{polys} = 65{,}536$
Metal is $1.48\times$ faster ($0.93$ \textmu s vs $1.38$ \textmu s).
\item WGSL via wgpu-native is faster than Metal at small batches
($n \in [16, 2688]$) because the wgpu-native command pipeline
setup is lighter than per-call metallib reload. WGSL never crosses
CPU on this hardware though — peaks at $0.89\times$ at
$n = 16{,}384$ — because per-thread arithmetic in the WGSL kernel
is uint32-emulated (WebGPU has no native uint64) and the constant
factor exceeds Metal's native uint64 path.
\item The Pulsar production hot path dispatches $M \cdot N_\text{vec}
\cdot \bar{D} = 8 \cdot 7 \cdot 48 = 2{,}688$ NTTs per signature. At
this batch size CPU and Metal are within $2\times$ of each other —
neither path is dramatically better.
\end{itemize}

\paragraph{What this paper claims.} Three empirical scaling laws
that should govern any honest claim about GPU acceleration in a
blockchain context. \S\ref{sec:scaling-laws} states them; they fall
out of the data with one fitted constant each.

\paragraph{What this paper does NOT claim.} CUDA throughput on
NVIDIA hardware is out of scope; we have no NVIDIA device on the
measurement host. The CUDA driver compiles cleanly and ships with
the same KAT contract, but absolute numbers will land in a follow-on
paper once a CUDA host is in CI. (See \S\ref{sec:honesty}.)
